\chapter{Mainstream scenarios: lifecycle과 이름의 미래}
\label{STC-S11}

미래를 하나의 확률 숫자로 압축하면 workload 차이를 잃는다. Consumer assistant는 external-memory
dominant이면서, enterprise coding agent는 hybrid promotion을 쓰고, foundation-model vendor는
periodic global refresh를 유지할 수 있다. 따라서 다섯 시나리오는 상호 배타적 시장 예측이 아니라
기술·경제 조건에 따른 system shape다.

더 중요한 언어적 주의가 있다. 이 lifecycle이 주류가 되어도 ``sleep-time compute''이라는 명칭이
승리할 필요는 없다. 공개 system은 dreaming, reflection, synthesis, memory, self-evolution,
consolidation이라는 서로 다른 말을 쓴다.\claim{STC-C042}
\autocite{srcstc0052,srcstc0054,srcstc0060,srcstc0063} 따라서 adoption은 keyword count가 아니라
deferred execution, wake-derived input, durable destination, later-wake reuse라는 operational signature로
추적해야 한다.

\section{다섯 개의 조건부 시나리오}

\begin{longtable}{L{0.15\textwidth} L{0.23\textwidth} L{0.25\textwidth} L{0.22\textwidth}}
\caption{시나리오별 trigger, leading indicator, falsifier}\label{tab:mainstream-scenarios}\\
\toprule
Scenario & Trigger & Leading indicators & Falsifiers \\
\midrule
\endfirsthead
\toprule
Scenario & Trigger & Leading indicators & Falsifiers \\
\midrule
\endhead
External dominant & provenance/delete 요구, frequent base upgrade, retrieval 개선 & pattern/supersede/cadence API, portable temporal memory, memory incident tooling & procedural skill을 지속적으로 잃거나 retrieval cost가 고정 병목 \\
Hybrid promotion & calibrated reuse predictor, cheap adapter validation/routing, source lineage & promote/demote/rebuild API, adapter catalog, reuse sweep benchmark & dual-state consistency와 version churn이 saving을 상쇄 \\
Periodic refresh & global amortization 우세, local history retention 제한 & release cadence 단축, central synthetic/eval factory, source-based adapter regeneration & local procedure가 release 사이에 빠르게 변함 \\
STC niche & 높은 reuse가 일부 domain에만 존재, validation cost 큼 & domain/cohort adapter, ROI/pressure trigger, repeated-task gain & general assistant에서도 background synthesis가 표준화 \\
STC broad & hundreds-cycle 안정성, bounded delete/rollback, stable queue, migration economics & wake/sleep quota와 SLO, durable state handle, lifecycle scaling curves & validation·state fan-out·privacy cost가 benefit보다 빠르게 증가 \\
\bottomrule
\end{longtable}

\section{현재 base case}

동결된 증거와 가장 일치하는 조합은 세 줄이다.

\begin{enumerate}
  \item \textbf{Near term: external-memory dominant.} Background synthesis, dedup, conflict resolution,
  temporal graph, provenance가 foundation weight를 바꾸지 않고 먼저 확산한다.
  \item \textbf{Evolution: high-reuse stable skill에서 hybrid promotion.} External source-of-record를
  버리지 않고 adapter, expert, latent cache를 제한적으로 compile한다.
  \item \textbf{Coexistence: global common knowledge는 periodic refresh.} 모든 durable learning을
  사용자별 sleep job으로 분해하지 않는다.
\end{enumerate}

STC niche의 가장 강한 형태, 즉 general assistant에는 background transformation이 거의 없을 것이라는
예측은 OpenAI와 Mem0 evidence로 이미 약해졌다. 반대로 continuous per-user parametric learning이
표준이 될 것이라는 broad scenario는 repeated-cycle E3와 production E4가 없다.

\section{분기점을 미리 보여 주는 signal}

\begin{table}[htbp]
\centering
\caption{Quarterly signal dashboard}
\label{tab:scenario-signals}
\begin{tabularx}{\textwidth}{L{0.32\textwidth} Y L{0.25\textwidth}}
\toprule
관측할 signal & 의미 & 어느 방향을 강화하는가 \\
\midrule
public per-user weight/adapter update와 rollback trace & parametric lifecycle 성숙도 & hybrid/broad \\
memory provenance·supersede·delete API 확장 & external governance maturity & external/hybrid \\
reuse predictor calibration과 top-$k$ yield & compile economics & hybrid/niche \\
base upgrade 후 adapter/index rebuild cost & artifact lifetime & refresh 대 hybrid \\
100--1,000 cycle retention/plasticity curve & lifelong claim 가능성 & broad 또는 반증 \\
background job bytes, energy, queue age & 실제 infra workload & 모든 시나리오의 device 규모 \\
learned writer/router의 cross-domain result & main-model update 대안 & external/hybrid \\
on-device idle training trace & privacy-local fit & niche/broad \\
\bottomrule
\end{tabularx}
\end{table}

\section{시나리오에 강건한 연구 전략}

예측에 베팅하기보다 공통 분모를 먼저 만든다. Immutable wake event와 lineage, versioned candidate,
validation attestation, atomic publish/rollback, total-state accounting은 external-only에서도 필요하고
hybrid와 broad adoption에서는 더 중요하다. Algorithm research는 reuse·volatility·capacity sweep으로
어떤 item을 external에 남기고 무엇을 promote할지 학습한다. Hardware research는 특정 ``sleep
accelerator'' 명칭보다 state locality, mixed rewrite, snapshot, metadata tail latency를 측정한다.

\begin{cautionbox}[Scenario를 달력 예측으로 읽지 않기]
이 장은 2028년 market share를 예측하지 않는다. 각 시나리오가 성립하려면 어떤 공개 evidence가
나타나야 하고 무엇이 confidence를 낮추는지 미리 고정한다. 새 정보가 들어오면 narrative가 아니라
trigger와 falsifier를 갱신한다.
\end{cautionbox}
